% filepath: /home/kluo/Documents/repos/kluo/mathphy/homework/2025/hw1_v3.tex
\documentclass[11pt]{ctexart}

\usepackage[margin=1in]{geometry} 
\usepackage{amsmath,amsthm,amssymb,mathtools}
\usepackage{graphicx, multicol, array}
\usepackage{ctex} 
\usepackage{fancyhdr}
\usepackage{fontspec}
\usepackage{fourier-orns}
\usepackage{xcolor}
\usepackage{enumitem}
\usepackage[most]{tcolorbox}

% Fix headheight warning
\setlength{\headheight}{15.12pt}
\addtolength{\topmargin}{-3.12pt}

% Chinese punctuation handling
\let\oldjuhao=。
\catcode`\。=\active
\newcommand{。}{\ifmmode\text{\oldjuhao}\else\oldjuhao\fi}

\let\olddouhao=，
\catcode`\，=\active
\newcommand{，}{\ifmmode\text{\olddouhao}\else\olddouhao\fi}

% Mathematical symbols
\newcommand{\N}{\mathbb{N}}
\newcommand{\Z}{\mathbb{Z}}
\newcommand{\R}{\mathbb{R}}
\newcommand{\C}{\mathbb{C}}
\newcommand{\ii}{\mathrm{i}}

% Colors for enhanced visual hierarchy
\definecolor{headercolor}{RGB}{25,25,112}
\definecolor{problemcolor}{RGB}{70,130,180}
\definecolor{sectioncolor}{RGB}{47,79,79}
\definecolor{accentcolor}{RGB}{220,20,60}

% Enhanced header rule with ornaments
\renewcommand{\headrule}{%
\vspace{-8pt}\hrulefill
\raisebox{-2.0pt}{\quad\color{headercolor}\decofourleft\decotwo\decofourright\quad}\hrulefill}

% Date configuration
\newcommand{\duedate}{\zhdate{2025/9/22}}

% Enhanced problem box styling
\tcbset{
    problembox/.style={
        enhanced,
        colback=blue!3!white,
        colframe=problemcolor,
        colbacktitle=problemcolor,
        fonttitle=\bfseries\large,
        coltitle=white,
        boxrule=1.2pt,
        arc=3pt,
        left=8pt, right=8pt, top=8pt, bottom=8pt,
        toptitle=6pt, bottomtitle=6pt,
        attach boxed title to top left={yshift=-2mm,xshift=3mm},
        boxed title style={arc=2pt,boxrule=0.8pt}
    }
}

% Custom problem environment
\newtcolorbox{hwproblem}[1]{
    problembox,
    title=题目 #1
}

\pagestyle{fancy}

% Header and footer setup with enhanced styling
\fancyhf{} % clear all fields
\fancyhead[L]{\color{headercolor}\textbf{数学物理方法}}
\zhnumsetup{style={Traditional,Financial}}
\fancyhead[R]{\color{headercolor}\textbf{课后习题\zhnumber{1}}}

\fancyfoot[L]{\small\color{gray}截止日期: \duedate}
\fancyfoot[C]{\color{headercolor}\textbf{\thepage}}
\fancyfoot[R]{\small\color{gray}任课教师: 罗凯}

% Enhanced title page styling
\title{
    \vspace{-1.5cm}
    {\Huge\color{headercolor}\textbf{数学物理方法}}\\
    \vspace{0.3cm}
    {\color{accentcolor}\rule{8cm}{1pt}}\\
    \vspace{0.3cm}
    {\Large\color{sectioncolor}课后习题一}
}
\author{
    \begin{tabular}{rl}
        \textbf{\color{sectioncolor}任课教师:} & 罗凯\\[0.2cm]
        \textbf{\color{sectioncolor}截止日期:} & \duedate
    \end{tabular}
}
\date{}

\begin{document}

% Custom enumerate styles with better spacing
\setlist[enumerate,1]{label=\color{problemcolor}(\arabic*), leftmargin=*, itemsep=0.6em, parsep=0.2em}
\setlist[enumerate,2]{label=\color{problemcolor}(\arabic{enumi}.\arabic*), leftmargin=*, itemsep=0.4em, parsep=0.1em}
\setlist[itemize]{label=\color{problemcolor}$\bullet$, leftmargin=*, itemsep=0.4em, parsep=0.1em}

\maketitle
\thispagestyle{fancy}

\begin{hwproblem}{1}
写出下列复数的\textbf{实部}、\textbf{虚部}、\textbf{模}和\textbf{辐角}。
\begin{enumerate}
    \item $1 + \ii \sqrt{3}$
    \item $\displaystyle\frac{1-\ii}{1+\ii}$
    \item $e^{z}$（其中 $z = x + \ii y$）
    \item $e^{1+\ii}$
    \item $e^{\ii \varphi(x)}$，其中 $\varphi(x)$ 是实数 $x$ 的实函数
\end{enumerate}
\end{hwproblem}

\vspace{0.5cm}

\begin{hwproblem}{2}
证明以下\textbf{代数运算规律}：
\begin{enumerate}
    \item 复数的加法和乘法满足\textbf{结合律}，即
    \begin{itemize}
        \item $(z_1 + z_2 ) + z_3 = z_1 +( z_2 + z_3)$
        \item $(z_1 \cdot z_2) \cdot z_3 = z_1 \cdot (z_2 \cdot z_3)$
    \end{itemize}
    \item 复数乘积的共轭等于复数共轭的乘积，即
    \[(z_1 \cdot z_2)^* = z_1^* z_2^*\]
\end{enumerate}
\end{hwproblem}

\vspace{0.5cm}

\begin{hwproblem}{3}
找出复平面上满足方程
\[|z - \ii a| = \lambda |z + \ii a|, \quad \lambda > 0\]
的所有点 $z = x + \ii y$。请分 $\lambda$ 的\textbf{三种情况}分别绘制图形：
\begin{enumerate}
    \item $\lambda < 1$（阿波罗尼乌斯圆，内部情况）
    \item $\lambda > 1$（阿波罗尼乌斯圆，外部情况）
    \item $\lambda = 1$（垂直平分线）
\end{enumerate}
\end{hwproblem}

\vspace{0.5cm}

\begin{hwproblem}{4}
若 $|z| = 1$，$a, b$ 为任意复数，试证明
\[\left| \frac{az + b}{\bar{b}z + \bar{a}} \right| = 1\]
\textit{提示：考虑使用复数模的性质和共轭运算。}
\end{hwproblem}

\vspace{0.5cm}

\begin{hwproblem}{5}
试用 $\cos\varphi$、$\sin\varphi$ 表示 $\cos 4\varphi$。

\textit{提示：利用欧拉公式 $e^{\ii\varphi} = \cos\varphi + \ii\sin\varphi$ 和二项定理。}
\end{hwproblem}

\vspace{0.5cm}

\begin{hwproblem}{6}
求下列方程的\textbf{所有根}，并在复平面上画出它们的位置。
\begin{enumerate}
    \item $z^4 + 1 = 0$
    \item $z^2 + 2z\cos\lambda + 1 = 0$，其中 $0 < \lambda < \pi$
\end{enumerate}
\textit{要求：用极坐标形式表示根，并标注在复平面图中。}
\end{hwproblem}

\vspace{0.5cm}

\begin{hwproblem}{7}
验证反正切函数的复数表示：
\[\tan^{-1}(z) = \frac{\ii}{2}[\ln(1-\ii z) - \ln(1+\ii z)]\]
\textit{提示：从 $\tan w = z$ 出发，利用欧拉公式推导。}
\end{hwproblem}

\vspace{0.5cm}

\begin{hwproblem}{8}
试证明\textbf{极坐标形式}的柯西-黎曼方程：
\[\left\{\begin{aligned}
\frac{\partial u}{\partial \rho} &= \frac{1}{\rho} \frac{\partial v}{\partial \varphi}\\[0.3em]
\frac{1}{\rho} \frac{\partial u}{\partial \varphi} &= -\frac{\partial v}{\partial \rho}
\end{aligned}\right.\]
其中 $z = \rho e^{\ii\varphi}$，$f(z) = u(\rho,\varphi) + \ii v(\rho,\varphi)$。
\end{hwproblem}

\vspace{0.5cm}

\begin{hwproblem}{9}
已知解析函数 $f(z)$ 的实部 $u(x, y)$ 或虚部 $v(x, y)$，求该\textbf{解析函数}。
\begin{enumerate}
    \item $u = e^x \sin y$
    \item $u = x^2 - y^2 + xy$，且 $f(0) = 0$
\end{enumerate}
\textit{提示：利用柯西-黎曼条件和调和函数的性质。}
\end{hwproblem}

\vfill

% Enhanced footer section
\begin{center}
\begin{tcolorbox}[
    colback=gray!5!white,
    colframe=gray!50!black,
    boxrule=0.5pt,
    arc=2pt,
    width=0.9\textwidth,
    left=10pt, right=10pt, top=5pt, bottom=5pt
]
\centering
\textcolor{sectioncolor}{\large\textbf{作业要求}}\\[0.3cm]
\begin{itemize}[label=\color{accentcolor}$\star$, leftmargin=2em]
    \item 请用\textbf{工整清楚}的字迹书写，\textbf{步骤完整}，逻辑清晰
    \item 复数计算请给出\textbf{详细过程}，图形绘制要\textbf{准确标注}
    \item 证明题需要\textbf{严格的数学推理}，避免跳步
    \item 如有疑问，请及时联系任课教师
\end{itemize}
\end{tcolorbox}
\end{center}

\end{document}